% QCM
Soit l'équation de la réaction~:
\[
    \ch{C3H8_{(g)}+5 O2_{(g)} -> 3 CO2_{(g)} + 4 H2O_{($\ell$)}}
\]

Les quantités initiales de réactifs sont pour le propane $(\ch{C3H8})$
$\qty{0.2}{\mol}$ et pour le dioxygène $\qty{0.8}{\mol}$.

On note $3x~\unit{\mol}$ la quantité de dioxyde de carbone formé (où $x$ est
l'avancement de la réaction)~:
\begin{enumerate}
    \item La quantité de dioxygène consommé est~: \par
          \begin{itemize*}[label=$\square$,itemjoin=\qquad]
              \item $5x$~; \item $x/5$~; \item $x$.
          \end{itemize*}
    \item La quantité de propane restant est~: \par
          \begin{itemize*}[label=$\square$,itemjoin=\qquad]
              \item $x$~; \item $0,2-x$~; \item $3 x$.
          \end{itemize*}
    \item Le réactif limitant est~: \par
          \begin{itemize*}[label=$\square$,itemjoin=\qquad]
              \item le propane~; \item le dioxygène.
          \end{itemize*}
    \item L'avancement maximal vaut~: \par
          \begin{itemize*}[label=$\square$,itemjoin=\qquad]
              \item $\qty{0.10}{\mol}$~;
              \item $\qty{0.16}{\mol}$~;
              \item $\qty{0.2}{\mol}$.
          \end{itemize*}
\end{enumerate}


